\section{Giới thiệu}

\subsection{Bối cảnh và mục tiêu}

Bài tập lớn 2 kế thừa kết quả từ BTL1 (44 mệnh đề đã được tách và phân tích cú pháp) để thực hiện các tác vụ \textbf{trích xuất thông tin ngữ nghĩa cấp cao}. Mục tiêu là chuyển đổi các mệnh đề hợp đồng thành các biểu diễn có cấu trúc, nắm bắt \textit{``ai làm gì với ai''} và \textit{chức năng pháp lý} của từng mệnh đề.

Ba tác vụ chính:
\begin{enumerate}
    \item \textbf{Custom NER (2.1)}: Nhận dạng và phân loại thực thể chuyên ngành pháp lý (bên ký, số tiền, ngày tháng, tỷ lệ, điều khoản phạt, luật áp dụng).
    \item \textbf{Semantic Role Labeling -- SRL (2.2)}: Gán vai nghĩa cho các đối số của từng vị từ trong mệnh đề (Agent, Theme, Recipient, Time, Condition...).
    \item \textbf{Clause Intent Classification (2.3)}: Phân loại intent pháp lý của mệnh đề (Obligation, Prohibition, Right, Termination Condition).
\end{enumerate}

\subsection{Dữ liệu}

\begin{table}[H]
    \centering
    \caption{Tổng quan dataset sử dụng trong BTL2}
    \label{tab:dataset_overview}
    \begin{tabular}{|l|l|r|l|}
        \hline
        \textbf{Tác vụ} & \textbf{Nguồn} & \textbf{Số mẫu} & \textbf{File} \\ \hline
        NER & CUAD + LLM synthetic & 308 & \texttt{ner\_training\_data.json} \\ \hline
        Intent & CUAD dataset & $\sim$480 & \texttt{intent\_training\_data.json} \\ \hline
        SRL & Pretrained (HuggingFace) & -- & Không cần training data \\ \hline
    \end{tabular}
\end{table}

\subsection{Mô hình và thư viện}

\begin{itemize}
    \item \textbf{NER}: Fine-tune \texttt{dslim/bert-base-NER} (HuggingFace Transformers + PyTorch).
    \item \textbf{SRL}: Pretrained \texttt{yeomtong/srl\_bert\_model} (BERT-base-cased).
    \item \textbf{Intent -- Baseline}: TF-IDF + Logistic Regression (scikit-learn).
    \item \textbf{Intent -- Nâng cao}: Fine-tune \texttt{distilbert-base-uncased}.
\end{itemize}